\section{Triple-Cert Adoption}
\label{sec:triple-cert}

\Zoo{} 3.0 inherits the \Lux{} \Quasar{} 3.0 triple-cert lane defined
by LP-020 \cite{lp-020}, with each cert layer specified in its own LP:
LP-075 for $\BLS{}_{12\text{-}381}$ \cite{lp-075}, LP-073 for
\Ringtail{} \cite{lp-073}, LP-070 for $\MLDSA{}\text{-}65$
\cite{lp-070}.

\subsection{The three lanes}

A \Quasar{} 3.0 cert is the conjunction of three independent proofs:

\begin{enumerate}[leftmargin=1.4em]
  \item \textbf{$\BLS{}_{12\text{-}381}$ aggregate (LP-075).}
        Classical fast-path proof. 48 bytes. Pairing-based.
        Falsifies under a CRQC; remains useful as the low-latency
        check that gates pre-final block propagation.
  \item \textbf{\Ringtail{} threshold (LP-073).}
        Post-quantum threshold proof built on Ring-LWE. Threshold
        $t$-of-$n$ aggregation across the validator set. Approximately
        2.0--2.4 KB depending on threshold and committee size.
  \item \textbf{$\MLDSA{}\text{-}65$ identity (LP-070).}
        Per-validator post-quantum identity proof, FIPS 204 conforming.
        Approximately 3.3 KB. Establishes which validator authorised
        a given cert layer.
\end{enumerate}

The triple-cert is the conjunction of all three. A block with only
$\BLS{}$ is invalid in 3.0 even if a majority of stake signed it
classically.

\subsection{Why triple, not double}

A double cert (classical + one PQ) leaves a single-point failure on
the post-quantum side. If the chosen lattice scheme has a structural
flaw that downgrades its security, the chain has no second PQ proof
to fall back on. \Quasar{} 3.0's design intent is to require two
independent PQ assumptions: \Ringtail{} relies on Ring-LWE hardness;
$\MLDSA{}$ relies on the hardness of M-LWE and M-SIS. A simultaneous
break of both --- on different lattice problems with different
parameter regimes --- is materially less likely than either
individually. The classical $\BLS{}$ layer remains for fast-path
liveness while the chain is healthy, and as a defence-in-depth
witness while a CRQC is plausible but not yet operational.

\subsection{What 3.0 inherits unchanged}

\Zoo{} 3.0 does not modify \Quasar{} 3.0 itself; \Zoo{} consumes it
as a chain-level service. The pieces \Zoo{} inherits from LP-020:

\begin{itemize}[leftmargin=1.2em]
  \item \textbf{TripleSignRound1.} The kernel that orchestrates parallel
        execution of the three cert lanes within a single round.
  \item \textbf{Photon committee selection (LP-105).} VRF-driven
        selection of the round's committee, deterministic given
        the previous \Quasar{}Cert.
  \item \textbf{Wave threshold voting (LP-105).} Threshold accumulation
        of signatures inside each lane.
  \item \textbf{Focus confidence accumulation (LP-105).} The check that
        all three lanes have crossed their respective thresholds before
        the cert is finalised.
\end{itemize}

\subsection{Where Zoo plugs in}

Zoo's L2 chain settles state commitments to the L1 \Quasar{}Cert. In
3.0, those commitments are themselves signed under the triple-cert
discipline, so a \Zoo{} block is final when:

\begin{enumerate}[leftmargin=1.4em]
  \item the \Zoo{} sequencer signs the block with $\BLS{}$, and
  \item the \Zoo{} sequencer submits the block hash to \Lux{}
        \Quasar{} 3.0 for inclusion, and
  \item a quorum of \Lux{} validators emit a \Quasar{}Cert that
        anchors the \Zoo{} block.
\end{enumerate}

The L2 inherits the L1's PQ guarantees through the cert. Pre-3.0,
\Zoo{} L2 also signed the inclusion proof with classical-only
multisig at the operator boundary; 3.0 adds the requirement that
the operator multisig be threshold $\MLDSA{}\text{-}65$ rather than
ECDSA.

\subsection{Verifier interface}

A \Zoo{} 3.0 light client must verify three proofs per epoch:

\begin{lstlisting}[language=C++,caption={Triple-cert verifier sketch}]
struct QuasarCert3 {
    bls::Aggregate           bls_proof;       // 48 bytes
    ringtail::ThresholdSig   ringtail_proof; // ~2.0-2.4 KB
    mldsa::Aggregate         mldsa_proof;    // ~3.3 KB * t
    Hash                     block_root;
};

bool verify_3_0(const QuasarCert3& c, const ValidatorSet& vs) {
    if (!bls_verify(c.bls_proof, c.block_root, vs))      return false;
    if (!ringtail_verify(c.ringtail_proof, c.block_root, vs)) return false;
    if (!mldsa_verify(c.mldsa_proof, c.block_root, vs))  return false;
    return true;
}
\end{lstlisting}

A failure on any lane fails the cert. There is no ``best of'' fallback.

\subsection{Performance budget}

The dominant cost in cert verification under 3.0 is the
$\MLDSA{}\text{-}65$ aggregate, because it scales linearly in the
number of distinct signers in the threshold whereas $\BLS{}$ aggregates
to a single 48-byte object. Empirically, on the \Lux{} validator
hardware class targeted for 3.0, the verifier closes at well below
the per-block latency budget (\cite{lp-105}, \cite{lp-020}). \Zoo{}'s
client side picks up the cost as a one-time verification per finalised
\Lux{} epoch; per-transaction cost is unchanged.
